\documentclass[12pt,a4paper]{article}

\usepackage{ctex}
\usepackage{amsmath,amssymb,amsfonts}
\usepackage{graphicx}
\usepackage{booktabs}
\usepackage{geometry}
\usepackage{setspace}
\usepackage{hyperref}
{{MAJOR_PACKAGES}}

\geometry{a4paper,left=3.17cm,right=3.17cm,top=2.54cm,bottom=2.54cm}
\onehalfspacing

\title{{{TITLE}}}
\author{{{AUTHOR}}}
\date{{{DATE}}}

\begin{document}

\maketitle

\begin{abstract}
这里写摘要内容（300--500 字）。简要说明研究背景、方法、结果与结论。

\noindent\textbf{关键词：}关键词1；关键词2；关键词3
\end{abstract}

\tableofcontents
\newpage

\section{绪论}

\subsection{研究背景}
在这里写研究背景...

\subsection{研究意义}
在这里写研究意义...

\section{相关工作}
在这里综述国内外研究现状...

\section{研究方法}
在这里描述本文的方法...

\section{实验结果与分析}

\subsection{实验设置}
在这里描述数据集、评价指标与实验环境...

\subsection{结果分析}
在这里分析实验结果，可插入三线表：

\begin{table}[htbp]
\centering
\begin{tabular}{lcc}
\toprule
方法 & 准确率 & 召回率 \\
\midrule
方法A & 0.912 & 0.887 \\
方法B & 0.934 & 0.901 \\
\bottomrule
\end{tabular}
\caption{不同方法的性能对比}
\label{tab:comparison}
\end{table}

\section{结论与展望}
在这里总结全文并展望未来工作...

\begin{thebibliography}{99}
\bibitem{ref1} 作者. 文献标题[J]. 期刊名, 年份, 卷(期): 页码.
\bibitem{ref2} 作者. 文献标题[M]. 出版社, 年份.
\end{thebibliography}

\section*{致谢}
在这里写致谢...

\end{document}
